%!TEX root = ../../main.tex
% ============================================================
% 统计图表 / 正比例图像
% 本文件由 main.tex 通过 \input 引入，请勿单独编译
% 重构日期: 2026-04-11
% ============================================================

\subsection{解答：探究总价与长度的正比例关系图像}

\vspace{0.6cm}

% 表格部分
\noindent\textbf{2. 一种花布每米售价 $8$ 元，购买 $2$ 米、$3$ 米$\cdots\cdots$分别需要多少元？}

\vspace{0.3cm}
\begin{center}
\renewcommand{\arraystretch}{1.5}
% 用较粗大一些的表格清晰展示
\begin{tabular}{|c|c|c|c|c|c|c|c|c|c|}
\hline
长度/米 & $1$ & $2$ & $3$ & $4$ & $5$ & $6$ & $7$ & $8$ & $\cdots$ \\
\hline
总价/元 & $8$ 
  & \textcolor{blue!80!black}{\textbf{16}}
  & \textcolor{blue!80!black}{\textbf{24}}
  & \textcolor{blue!80!black}{\textbf{32}}
  & \textcolor{blue!80!black}{\textbf{40}}
  & \textcolor{blue!80!black}{\textbf{48}}
  & \textcolor{blue!80!black}{\textbf{56}}
  & \textcolor{blue!80!black}{\textbf{64}}
  & $\cdots$ \\
\hline
\end{tabular}
\end{center}

\vspace{0.6cm}

% 左右分栏布局：左边题目与解答，右边放 TikZ 正比例函数图
\noindent
\begin{minipage}[t]{0.50\textwidth}
\vspace{0pt} % 贴顶对齐
\normalsize
（1）把上面的表格填写完整。\\
（2）根据表中数据，先在右图中描出长度和总价所对应的点，再把这些点依次连起来。\\
（3）购买这种花布的总价和长度成正比例吗？为什么？

\vspace{0.4cm}
\textcolor{blue!80!black}{\textbf{答：成正比例，因为总价随着长度的变化而变化，且两者的比值一定。}}
\end{minipage}
\hfill
\begin{minipage}[t]{0.45\textwidth}
\vspace{0pt}
\begin{center}
\begin{tikzpicture}[>=Stealth, scale=0.6] % 缩放适配版面右侧

  % 1. 绘制浅蓝色背景网格 (横向9格到x=9，纵向9格到y=9即代表72)
  \draw[step=1.0cm, cyan!50, line width=0.6pt] (0, 0) grid (9.0, 9.0);
  
  % 2. 绘制坐标轴 (黑色带箭头段略长出网格)
  % X 轴
  \draw[->, black, line width=1.0pt] (0, 0) -- (9.8, 0) node[right, font=\small] {长度/米};
  % Y 轴
  \draw[->, black, line width=1.0pt] (0, 0) -- (0, 9.6) node[right, font=\small, xshift=0.1cm] {总价/元};
  
  % 原点刻度 O (0)
  \node[below left, font=\footnotesize] at (0, 0) {$0$};

  % 3. 绘制坐标系横纵轴的刻度文字
  % 横轴：1 到 9
  \foreach \x in {1, 2, ..., 9} {
    \node[below, font=\footnotesize, yshift=-0.1cm] at (\x, 0) {$\x$};
  }
  
  % 纵轴：原题纵轴上对应0-9各格子，1代表8，9代表72
  \foreach \y/\val in {1/8, 2/16, 3/24, 4/32, 5/40, 6/48, 7/56, 8/64, 9/72} {
    \node[left, font=\footnotesize, xshift=-0.1cm] at (0, \y) {\val};
  }

  % 4. 描点连线作图 (答案部分：蓝色折线穿过各个交点)
  \draw[blue!80!black, line width=1.2pt] (1, 1) -- (8, 8);
  
  % 画出实心蓝点即为对应的交点
  \foreach \i in {1, 2, 3, 4, 5, 6, 7, 8} {
    \fill[blue!80!black] (\i, \i) circle (4pt);
  }

\end{tikzpicture}
\end{center}
\end{minipage}

\vspace{0.8cm}
\hrule
\vspace{0.4cm}

{\small\color{gray}
\textbf{解题说明：}\\
1. \textbf{填写表格：}已知花布单价为恒定的 $8\,\text{元/米}$。根据“总价 $=$ 单价 $\times$ 长度”，可以直接依次推算出购买长度 $2, 3, \cdots, 8$ 米所对应的总价分别为 $16, 24, \cdots, 64$ 元。填入表格内标蓝处即可。\\
2. \textbf{描点连线作图：}观察原版题目预设的空白坐标系基座，我们发现其十字网格刻度分配是：横轴单位格代表长度 $1\text{m}$，纵轴单位格恰好代表总价 $8\text{元}$。因此我们将表格上侧列作为横坐标值 $x$，下行对应列金额作为纵轴真实长度换算高度变量 $y$，得到的有序数对 $(1,8)、(2,16) \cdots$ 会地直接落在背景细网格的各个对角线正交汇处。在此处用明显的实心小圆点深描，再施以直尺连线使其平滑成直线。\\
3. \textbf{正比例函数定性判定解答：}因为在几何图像表面上，这两者的关系被最终绘制成了一条必过原点坐标 $(0,0)$ 的斜向倾斜直线（正比例函数第一象限图像特征）；而在代数根基关系逻辑上：由于总价随着长度成倍递增变化趋势，且商“$\text{总价} \div \text{长度} = \text{单价}(8\text{元/米} > 0)$”，其比率值是一个不可动摇的极常数量，由此彻底证出和满足了判断其正比例属性核心定语：“总价比长度为常数一”。
}

\vspace{0.6cm}

{\small\color{blue!70!black}
\textbf{绘图基本原理与步骤：}\\
\textbf{1. 面向报表排印的矩阵级高分栏：}本题包含了长宽不一的表格以及经典的左试卷题干、右描边图系的混合双排版需求。先在顶级流用 \texttt{tabular} 配合 \texttt{\textbackslash textcolor\{blue\}} 高亮醒目突显表格算术推导区，强加阅读辅助反馈。其后排版利用纯正的 \texttt{\textbackslash begin\{minipage\}} 施展极宽版面左右撕裂切割场法：左侧横向切出占空 \texttt{0.50} 承装长文短句试炼题文，右侧空投 \texttt{0.45} 装载包含复杂极坐标矢量的几何画廊 \texttt{tikzpicture}。它保证了即便后期修改文字，这两片浮动版图也将如同不可相互侵犯的平行宇宙，左右并持且上下顶端齐平约束，呈现出超越回车空格敲击法的几何排版美学。\\
\textbf{2. 自动化步进递推底纹以及简化数轴标记流：}由于原本考题插画底背景有很致密显眼的一大片蔚然蓝海网格，为了防止手绘多条平行线的繁冗，底图一键祭出超级引擎函数 \texttt{\textbackslash draw[step=1.0cm, cyan!50] (0,0) grid (9,9)}，这种直接的网格发生器仅用半行代码就算准出纵横整整十九条长线段相交生成的八十一宫格铺装地基。而在极易写错位的人为数字标注端，放弃手工定位，直接采用自带泛型矩阵性质的 \texttt{\textbackslash foreach} 字典序列扫描。像是那段 \texttt{1/8, 2/16} 构成的映射串列，让纵轴上的标签渲染变成类似解包 Python 列表元组一样，一边抽取出 $y$ 坐标刻度的几何高度 $1,2,3...$，一边向屏幕喷涌抛出渲染出它的实际物理表代表含义 $8,16,24...$ 这不仅缩短了冗长代码，而且防呆容错性使得后续如果物价改为 $10\text{元}$ 也只消在序列库里一键微调全局自动焕发。
}

%% ==============================
%% 第15题：最短路径问题（饮马问题及模型应用）
%% ==============================
\newpage
